Internal corrosion governs the residual life of production piping in
mature oil fields. Integrity programs quantify its progression through
the corrosion rate $v$, expressed in mils per year [mpy], and classify
the severity against normative limits: rates up to 2 [mpy] permit
routine monitoring, rates between 2 [mpy] and 8 [mpy] require
intensified inspection, and rates above 8 [mpy] demand intervention.
Field measurement campaigns produce databases with two structural
difficulties. First, the observations concentrate in the low-severity
range, so the critical regimes remain under-represented
\cite{he2009imbalanced}. Second, the number of useful records rarely
exceeds a few hundred after quality control.

Regression models trained on such data minimize a global error and,
consequently, underestimate the scarce high-severity observations. The
resulting predictions bias inspection planning toward the benign
regimes. Flexible ensemble regressors reduce the global error
\cite{breiman2001random} but hide the reasoning behind each
prediction, which conflicts with the audit requirements of integrity
management \cite{rudin2019stop}; on reduced databases their high
parameter count also raises the risk of overfitting.

Fuzzy inference systems offer the language that integrity engineers
already use: severity classes, gradual transitions, and IF--THEN rules
\cite{zadeh1965fuzzy,mamdani1975experiment}. Their practical
difficulty lies in the construction of the rulebase and of the
membership functions, which traditionally depend on manual
elicitation \cite{guillaume2001designing}. Decision trees resolve this
difficulty: a tree trained on the normative classes yields an explicit
partition of the input space whose branches read as rules
\cite{breiman1984cart}.

This study articulates both elements into a single strategy. Binomial
decision trees, one per normative limit, provide the classification
backbone; the tree thresholds define interval membership functions
whose transition widths follow the training-data distribution; and the
resulting rulebase feeds a fuzzy inference system whose consequents
remain bounded to the normative ranges. A nested cross-validation
protocol confines every selection step to inner folds and eliminates
the optimistic bias of tune-and-refit procedures
\cite{varma2006bias,cawley2010overfitting}.

The contributions are: (i) a reproducible procedure that converts
threshold-wise decision trees into an auditable fuzzy estimator of the
corrosion rate; (ii) a distribution-based construction of interval
membership functions with transitions anchored at the tree thresholds;
(iii) a record-level leakage audit and an evaluation protocol that
together guarantee leakage-free estimates of generalization for the
complete chain; and (iv) a comparison on identical folds against
regression methods suited to reduced databases, including the
sensitivity to the severe class that drives integrity decisions.
